\scanpage{295}
$C$ cat. abélienne

$S$ anneau $\mathbf Z$-gradué comm. à degrés $\geq0$

$\mathcal M$ \inserted{$=\mathcal M(C,S)$} catégorie des $C$-$S$-modules gr. \uline{\unsure{diffs.}}

$\mathcal M_N$, $\mathcal M^N$ définis comme plus haut.

\diagram{\begin{tikzcd}\mathcal M_N\arrow[r,"i_N"]&\mathcal M\end{tikzcd}}

admet \deleted{foncteur} adjoint à droite

\diagram{\begin{tikzcd}\mathcal M\arrow[r,"\tau_N"]&\mathcal M_N\end{tikzcd}}
\[ i_N\tau_N(P)\subset P \]

\deleted{qui admet}

Construire adjoint à droite $\varphi_N$ de $\tau_N$ :
\[
\operatorname{Hom}_{\mathcal M_N}(\tau_NP,\underset{\in\mathcal M_N}{Q})\simeq
\operatorname{Hom}_{\mathcal M}(P,\varphi_NQ)\ ?
\]
\margin{\unsure{1}}

Il faut \unsure{donc} prouver que pour $Q$ fixé, le premier mb, comme foncteur en $P$, est repr. Or il est clair que $\tau_N$ commute aux $\varinjlim$ qui sont \unsure{représent.bles} dans $\mathcal M$ donc le foncteur du premier mb transforme ces $\varinjlim$ en $\varprojlim$. Cela implique donc ce qu'on veut lorsque $C$ \uline{ou} $C^o$ est une ``cat. de Groth''\ldots

Comment \inserted{alors} décrire les $(\varphi_NQ)_i\in\operatorname{Ob}C$

On a pour tout $M\in\operatorname{Ob}\mathcal M$
\[ M_i\simeq\underline{\operatorname{Hom}}^*(S,M)_i \]
donc
\[
\begin{aligned}
(\varphi_NQ)_i&=\underline{\operatorname{Hom}}^*(S,\varphi_NQ)_i
\simeq\underline{\operatorname{Hom}}^*(S[-i],\varphi_NQ)_0\\
&\overset{\unsure{\text{quid pro}}}{=}\underline{\operatorname{Hom}}^*(\tau_N(S[-i]),Q)_0\\
&=(\underline{\operatorname{Hom}}^*((\tau_{N-i}(S))[-i],Q)\deleted{\unsure{\tau_N}})_0
\end{aligned}
\]
\margin{296}
